\documentclass[11pt]{amsart}
\usepackage{amsmath,amssymb,amsthm,booktabs}
\usepackage[margin=1.15in]{geometry}
\usepackage[hidelinks]{hyperref}
\newtheorem{theorem}{Theorem}[section]
\newtheorem{corollary}[theorem]{Corollary}
\newtheorem{proposition}[theorem]{Proposition}
\newtheorem{lemma}[theorem]{Lemma}
\theoremstyle{remark}
\newtheorem{remark}[theorem]{Remark}
\newcommand{\doi}[1]{\href{https://doi.org/#1}{doi:\nolinkurl{#1}}}
\newcommand{\eps}{\varepsilon}
\newcommand{\E}{\mathbb{E}}
\newcommand{\Z}{\mathbb{Z}}
\newcommand{\N}{\mathbb{N}}
\newcommand{\thit}{t_{\mathrm{hit}}}
\newcommand{\rhit}{\rho_{\mathrm{hit}}}
\setlength{\emergencystretch}{2.5em}

\begin{document}

\title[Hitting-time robustness of the discrete torus]{Hitting-time robustness of the discrete torus}
\author{Yitzchak Shmalo}
\address{Einstein Institute of Mathematics, The Hebrew University of Jerusalem, Jerusalem, Israel}
\email{yspennstate@gmail.com}
\subjclass[2020]{60J10, 60J27, 05C81}
\keywords{random walk, discrete torus, hitting time, adversarial perturbation, heat kernel, Markov chain robustness}

\begin{abstract}
Zuckerman asked how far an adversary who may change every transition probability of a random walk by a factor in
$[1-\eps,1+\eps]$, depending on the whole history, can raise its maximum expected hitting time. For the simple random
walk on the discrete torus $\Z_L^D$ with $L$ even, $D\ge3$ and $\eps=c/L$ we show that the strongest such adversary raises the
maximum expected hitting time by a factor tending to $((e^c-1)/c)^D$. An adversary that pushes every coordinate away from the
target independently is optimal up to a factor $1+o(1)$; under it, in continuous time, the distances of the coordinates from those of the target are independent copies of one birth--death chain on $\{0,\dots,L/2\}$, the folded chain. The proof that no
adversary does better by more than a factor $1+o(1)$ rests on a product inequality for the transition probabilities to $0$ of a birth--death chain on the half-line, proved
by coupling a killed time-inhomogeneous chain; at times of order at most $L^2/\log L$ it transfers to the folded chain up to
an error polynomially small in $L$. It follows that the hitting-time robustness of
$\Z_L^D$ is of order one over the diameter; the lower bound is the case of the torus of a conjecture of Zuckerman for all
graphs.
\end{abstract}

\maketitle

\section{Introduction}\label{sec:intro}

Let $G$ be a connected graph and $P$ the transition probabilities of the simple random walk on $G$. Following
Zuckerman~\cite{Zuc26}, a \emph{standard $\eps$-adversary} ($0<\eps<1$) moves the walk from its current vertex $v$ to a
neighbour $w$ with a probability of its choosing in $[(1-\eps)P(v,w),(1+\eps)P(v,w)]$, the probabilities from $v$
summing to one; the choice may depend on the whole history. Write $H_v$ for the hitting time of $v$, $h^+(u,v)$ for
the supremum of $\E_u H_v$ over all standard $\eps$-adversaries, $\thit^+(\eps)=\max_{u,v}h^+(u,v)$, and
$\thit=\max_{u,v}\E_uH_v$ for the unperturbed walk. The \emph{hitting-time robustness} is
\[
\rhit(C)=\sup\{\eps:\ \thit^+(\eps)\le C\,\thit\}.
\]
When $C$ is omitted, $C$ is a sufficiently large constant, as in~\cite[Definition~1.9]{Zuc26}.
Zuckerman proved $\rhit=\Omega(1/t_{\mathrm{mix}})$ for every graph~\cite[Theorem~1.10]{Zuc26}, where
$t_{\mathrm{mix}}$ is the mixing time (of the lazy walk when $G$ is bipartite~\cite[Remark~2.9]{Zuc26}), conjectured
that $\rhit=\Omega(1/\Delta)$ with $\Delta$ the diameter~\cite[Conjecture~1.16]{Zuc26}, proved this for trees~\cite[Theorem~1.11]{Zuc26}, and listed
grids among the natural graphs whose robustness is to be determined~\cite[Section~1.10]{Zuc26}. Counterexamples to the conjecture, on cubic graphs of logarithmic diameter and on regular graphs of diameter three, have been proposed in~\cite{Anon26}. On the torus
$\Z_L^D$ the mixing time is of order $L^2$ for fixed $D$~\cite{LPW17}, while $\Delta=DL/2$, so the general bound gives
$\rhit=\Omega(1/L^2)$ and leaves a gap of a factor $L$. On the other side, the adversary constructed on a regular graph in the proof
of~\cite[Theorem~4.5]{Zuc26}, a Metropolis chain whose holding probabilities are spread over the edges in proportion to
their probabilities, multiplies every transition probability by a factor in $[1-\eps,(1-\eps)^{-1}]$. For $\eps<1/2$
it is therefore a standard $\eps/(1-\eps)$-adversary, and under it the maximum expected hitting time is
$\Omega(n\,e^{\eps\bar\Delta})$, where $n$ is the number of vertices and $\bar\Delta$ the average distance. On
$\Z_L^D$ with $D\ge3$ we have $n=L^D$, $\bar\Delta$ of order $L$ and $\thit$ of order $L^D$, so this gives
$\rhit=O(1/L)$, and the gap is on the lower side. We close
this gap on $\Z_L^D$ with $L$ even and $D\ge3$, by identifying an adversary that is optimal up to a factor $1+o(1)$ and computing its
maximum expected hitting time. For the hypercube $\{0,1\}^d$ the strongest adversary is determined exactly in~\cite{ShmCube}: at $\eps=c/d$ it raises
the maximum expected hitting time by a factor tending to $2^c$.

We work on the torus $\Z_L^D$ with $L$ even and target $0$. Its folded coordinates are
$k_a=\min(x_a,L-x_a)\in\{0,\dots,m\}$, $m=L/2$. Consider the \emph{product adversary} $Q$ that in each coordinate puts the larger
weight $(1+\eps)/(2D)$ on the step that increases $k_a$ and the smaller weight $(1-\eps)/(2D)$ on the step that
decreases it, and leaves the weights $1/(2D)$ unchanged when $k_a\in\{0,m\}$, where the two steps in coordinate $a$
both increase or both decrease $k_a$. These weights lie in the allowed range and sum to one, so $Q$ is a standard
$\eps$-adversary, and it does not depend on the history. On the hypercube the supremum is attained by the adversary
that at every vertex sends the least mass toward the target that the bounds allow~\cite{ShmCube}. On the torus $Q$ is the analogous adversary, but at the sites where a coordinate is $0$ or $m$ it is not the best one-step choice against its own hitting times, and the
comparison described after Corollary~\ref{cor:rho} is needed. The condition $D\ge3$ makes the simple random walk on $\Z^D$ transient; for $D=2$ the order of $\thit$ is $L^2\log L$ (Section~\ref{sec:checks}). We take
$L$ even so that at the far end $k_a=m$ both steps in coordinate $a$ decrease $k_a$. Here and below $O$ hides constants that depend only on $c$ and $D$, the constants in
$\Omega$ and $\Theta$ may also depend on $C$, and $o(1)$ refers to $L\to\infty$.

\begin{theorem}\label{thm:main}
Let $D\ge3$ and $c>0$ be fixed and let $L\to\infty$ through even integers. For the simple random walk on $\Z_L^D$,
\[
\lim_{L\to\infty}\frac{\thit^+(c/L)}{\thit}=\Bigl(\frac{e^c-1}{c}\Bigr)^{D}.
\]
More precisely, $\thit^+(c/L)=\max_x\E^Q_xH_0\,\bigl(1+O((\log L)^{D/2}/L)\bigr)$, and
$L^{-D}\max_x\E^Q_xH_0\to G_D\bigl((e^c-1)/c\bigr)^D$, where $G_D=\int_0^\infty P_0(S_t=0)\,dt$ for the
continuous-time simple random walk $S$ on $\Z^D$ with jump rate one.
\end{theorem}

\begin{corollary}\label{cor:rho}
Under the hypotheses of Theorem~\ref{thm:main}, for every $C>1$, $L\,\rhit(C)\to r_D(C)$, the unique root of
$((e^r-1)/r)^D=C$. In particular $\rhit(C)=\Theta(1/\Delta)$ on $\Z_L^D$.
\end{corollary}

The corollary follows from the theorem because $\thit^+(\eps)$ is nondecreasing in $\eps$ (the set of adversaries
grows) and $r\mapsto((e^r-1)/r)^D$ is continuous and strictly increasing from $1$ to $\infty$ on $(0,\infty)$: for
$c_1<r_D(C)<c_2$ the theorem gives $\thit^+(c_1/L)<C\thit<\thit^+(c_2/L)$ for large $L$, hence
$c_1/L\le\rhit(C)\le c_2/L$.

Under $Q$ the folded coordinates, run in continuous time at rate $1/D$ each, are independent copies of the
\emph{folded chain} $X$, a birth--death chain on $\{0,\dots,m\}$ with drift $\eps$ away from $0$, and $V_Q(x)=\E^Q_xH_0$ can be computed; by symmetry it depends on $x$ only through
$k=(k_1,\dots,k_D)$, and we write $V_Q(k)$.
Let $B$ send a function $f$ to $Bf(x)$, the largest value of $\sum_yP'(y)f(y)$ over the one-step laws $P'$ an
adversary may use at $x$. A supermartingale argument (Section~\ref{sec:reduction}) shows that no adversary beats $Q$
by more than the factor $1/(1-\bar d)$, where $\bar d=\max_{x\ne0}(BV_Q(x)-QV_Q(x))$. The difficulty is that $QV_Q$ falls
short of $BV_Q$ at some sites. At a site where some coordinate equals $0$ or $m$, both steps in that coordinate change $k_a$ in the same
direction, and the adversary may gain by moving weight to another coordinate. The deficit $BV_Q(x)-QV_Q(x)$ is at most
$\eps/D$ times a sum of differences such as $V_Q(k+e_a)-V_Q(k+e_i)$, where $e_a$ is the $a$-th unit vector, $k_a=0$
and $1\le k_i\le m-1$: a step out of $0$ in a fresh coordinate against a further step out in an interior one.
Monotonicity of $V_Q$ in each coordinate does not compare the two. The new ingredient is an inequality for the chain
$X^\infty$ on $\{0,1,2,\dots\}$ that jumps from $0$ to $1$ at rate $1$ and otherwise up at rate $(1+\eps)/2$ and
down at rate $(1-\eps)/2$: with $f_t(x)=P_x(X^\infty_t=0)$,
\begin{equation}\label{eq:key}
f_t(1)\,f_t(j)\ \ge\ f_t(0)\,f_t(j+1)\qquad(t>0,\ j\ge1).
\end{equation}
At $\eps=0$, $f_t(j)=e^{-t}I_j(t)$ with $I_j$ the modified Bessel functions, and~\eqref{eq:key} follows by chaining the
Tur\'an-type inequality $I_\nu^2\ge I_{\nu-1}I_{\nu+1}$~\cite{BP}. For $\eps>0$ we prove it by a coupling of a
time-inhomogeneous chain (Lemma~\ref{lem:halfline}). With it, and with Delmotte's Gaussian
bounds~\cite{Del99} for the on-diagonal kernel of $X$, the deficit satisfies $\bar d=O((\log L)^{D/2}/L)$, so $Q$ is
optimal up to a factor $1+o(1)$, and Theorem~\ref{thm:main} reduces to computing $\max_xV_Q(x)$.

\emph{Outline.} Section~\ref{sec:reduction} proves the reduction to $Q$ and the form of the deficit. Section~\ref{sec:folded} collects facts about $X$, and Section~\ref{sec:halfline} proves~\eqref{eq:key}. Section~\ref{sec:deficit} uses it to bound $\bar d$ by $O((\log L)^{D/2}/L)$, and
Section~\ref{sec:value} computes
$\max V_Q$. Section~\ref{sec:checks} gives numerical values of $\thit^+(1/L)/\thit$ on $\Z_L^2$ and of $\max_xV_Q(x)/\thit$ on $\Z_L^3$. The dependencies are as follows.
\begin{center}
\begin{tabular}{lll}
\toprule
Section & Result & Uses\\ \midrule
\ref{sec:reduction} & reduction to $Q$, form of $d$ (Lemmas~\ref{lem:reduction}, \ref{lem:deficit}) & Lemma~\ref{lem:folded}(iii)\\
\ref{sec:folded} & the folded chain (Lemmas~\ref{lem:folded}--\ref{lem:kernel}) & \cite{Del99}, on-diagonal\\
\ref{sec:halfline} & product inequality~\eqref{eq:key} (Lemma~\ref{lem:halfline}) & --\\
\ref{sec:deficit} & $\bar d=O((\log L)^{D/2}/L)$ (Proposition~\ref{prop:deficit}) & Sections~\ref{sec:reduction}--\ref{sec:halfline}\\
\ref{sec:value} & $\max_xV_Q(x)$ (Proposition~\ref{prop:value}), proof of Theorem~\ref{thm:main} & Section~\ref{sec:folded}\\
\ref{sec:checks} & numerical values for $D=2,3$ & --\\
\bottomrule
\end{tabular}
\end{center}

\section{Reduction to the product adversary}\label{sec:reduction}

By transitivity it suffices to consider the target $0$. Write $V=V_Q$, $V(x)=\E^Q_xH_0$; by symmetry $V(x)$ depends
only on the folded coordinates $k=k(x)$ and is invariant under permutations of them, and we write $V(k)$. First-step
analysis gives $QV(x)=V(x)-1$ for $x\ne0$. For a function $f$ on $\Z_L^D$ let
\[
B f(x)=\max\Bigl\{\sum_yP'(y)f(y)\Bigr\},
\]
the maximum over probability vectors $P'$ on the $2D$ neighbours of $x$ with $P'(y)\in[(1-\eps)/(2D),(1+\eps)/(2D)]$,
and call $d(x)=BV(x)-QV(x)\ge0$, $x\ne0$, the \emph{deficit}.

\begin{lemma}\label{lem:reduction}
If $\bar d=\max_{x\ne0}d(x)<1$, then $\max_xV(x)\le\thit^+(\eps)\le\max_xV(x)/(1-\bar d)$.
\end{lemma}

\begin{proof}
The lower bound holds because $Q$ is a standard $\eps$-adversary. For the upper bound fix any standard
$\eps$-adversary, a start $x$, and let $\xi_t$ be the walk it produces and $\mathcal F_t$ its history. On $\{t<H_0\}$,
$\E[V(\xi_{t+1})\mid\mathcal F_t]\le BV(\xi_t)=V(\xi_t)-1+d(\xi_t)\le V(\xi_t)-(1-\bar d)$. Hence
$V(\xi_{t\wedge H_0})+(1-\bar d)(t\wedge H_0)$ is a supermartingale, so $(1-\bar d)\E_x[t\wedge H_0]\le V(x)$, and
monotone convergence gives $\E_xH_0\le V(x)/(1-\bar d)$.
\end{proof}

The deficit has a simple form. Fix $x\ne0$ with folded coordinates $k$ and split the coordinates into
$I=\{a:1\le k_a\le m-1\}$, $S_0=\{a:k_a=0\}$ and $S_m=\{a:k_a=m\}$, with $s_0=|S_0|$, $s_m=|S_m|$. For $a\in I$ the
two neighbours in direction $a$ have values ($e_a$ is the $a$-th unit vector) $o_a=V(k+e_a)$ (outward) and $w_a=V(k-e_a)$ (inward). For $a\in S_0$ both
neighbours have the value $A=V(k+e_a)$, which by permutation symmetry is the same for all $a\in S_0$; for $i\in S_m$
both have the value $M=V(k-e_i)$, the same for all $i\in S_m$. By Lemma~\ref{lem:folded}(iii) below,
\begin{equation}\label{eq:order}
w_a\le V(k)\le o_a\quad(a\in I),\qquad M\le V(k)\le A.
\end{equation}

\begin{lemma}\label{lem:deficit}
If $S_m=\emptyset$, then $d(x)\le(\eps/D)\,s_0\bigl(A-\min_{a\in I}o_a\bigr)^+$ (and $d(x)=0$ if also $s_0=0$).
In general $d(x)\le(\eps/D)\bigl(s_0\,(A-V(k))+s_m\,(V(k)-M)\bigr)$.
\end{lemma}

\begin{proof}
Write $P'(y)=(1+\eps\sigma_y)/(2D)$ with $\sigma_y\in[-1,1]$ and $\sum_y\sigma_y=0$. The maximum of the linear function
$\sum_y\sigma_yV(y)$ over this polytope is attained at $\sigma=+1$ on $D$ neighbours with the largest values and $-1$ on
the rest. Group the $2D$ neighbour values into $D$ pairs, one per coordinate: $(o_a,w_a)$ for $a\in I$, $(A,A)$ for
$a\in S_0$, $(M,M)$ for $a\in S_m$. Let $F$ be the multiset of first entries and $S$ that of second entries. Then $Q$
corresponds to $\sigma=+1$ on $F$ and $-1$ on $S$ (for the equal pairs this changes nothing), so with
$\Sigma_{\mathrm{top}}$ the sum of the $D$ largest values,
\[
d(x)=\frac{\eps}{2D}\Bigl(2\Sigma_{\mathrm{top}}-\textstyle\sum F-\sum S\Bigr)-\frac{\eps}{2D}\Bigl(\sum F-\sum S\Bigr)
=\frac{\eps}{D}\Bigl(\Sigma_{\mathrm{top}}-\textstyle\sum F\Bigr).
\]
Order $S$ decreasingly, $s_{(1)}\ge s_{(2)}\ge\cdots$, and $F$ increasingly, $f_{(1)}\le f_{(2)}\le\cdots$. A set of
$D$ largest values contains the $r$ largest elements of $S$ and misses the $r$ smallest of $F$ for some $r$, so
$\Sigma_{\mathrm{top}}-\sum F=\sum_{i\le r}(s_{(i)}-f_{(i)})\le\sum_i(s_{(i)}-f_{(i)})^+$. By~\eqref{eq:order} a
positive term has $s_{(i)}=A$ or $f_{(i)}=M$: the remaining elements of $S$ are $\le V(k)$ and those of $F$ are
$\ge V(k)$. There are at most $s_0$ terms with $s_{(i)}=A$ and at most $s_m$ with $f_{(i)}=M$. If $S_m=\emptyset$, each
positive term is $A-o_a$ for some $a$, which gives the first bound. In general a term $A-o_a$ is at most $A-V(k)$, a
term $w_a-M$ at most $V(k)-M$, and a term $A-M$ equals $(A-V(k))+(V(k)-M)$, which gives the second.
\end{proof}

So Theorem~\ref{thm:main} follows from two estimates, proved in Sections~\ref{sec:deficit} and~\ref{sec:value}:
\begin{itemize}
\item[(a)] $\bar d=O((\log L)^{D/2}/L)$;
\item[(b)] $L^{-D}\max_xV(x)\to G_D((e^c-1)/c)^D$.
\end{itemize}
The case $c=0$ of (b), where $Q$ is the simple random walk, gives $\thit=G_DL^D(1+o(1))$, and the theorem follows from
Lemma~\ref{lem:reduction}.

\section{The folded chain}\label{sec:folded}

Let $X$ be the continuous-time chain on $\{0,\dots,m\}$ that jumps from $0$ to $1$ at rate $1$, from $m$ to $m-1$ at
rate $1$, and from $1\le k\le m-1$ to $k+1$ at rate $\alpha=(1+\eps)/2$ and to $k-1$ at rate $\beta=(1-\eps)/2$. Write
$p_s(j,k)=P_j(X_s=k)$ and $\pi$ for its stationary law. In continuous time with jump rate one, the walk $\xi$ under $Q$ has
folded coordinates $X^{(1)}_{t/D},\dots,X^{(D)}_{t/D}$ with $X^{(a)}$ independent copies of $X$, and the expected
hitting time of $0$ is the same as for the discrete-time walk, since the jump chains agree. Thus
\begin{equation}\label{eq:product}
P^Q_x(\xi_t=0)=\prod_{a=1}^Dp_{t/D}(k_a,0),\qquad \pi_Q(0)=\pi(0)^D,
\end{equation}
where $\pi_Q$ is the stationary law of $Q$ on $\Z_L^D$. Write $\hat p_s(j)=p_s(j,0)/\pi(0)$.

The chain is reversible, with $\pi(k)=\pi(0)\,\beta^{-1}(\alpha/\beta)^{k-1}$ for $1\le k\le m-1$ and $\pi(m)=\alpha\,\pi(m-1)$. With
$\eps=c/L$ we have $(\alpha/\beta)^m\to e^{c}$, so
\begin{equation}\label{eq:pi0}
\pi(0)^{-1}=1+\frac1\beta\sum_{k=1}^{m-1}\Bigl(\frac\alpha\beta\Bigr)^{k-1}+O(1)=L\,\frac{e^c-1}{c}\,(1+o(1)),
\end{equation}
and $\pi(k)/\pi(0)\in[C^{-1},C]$ for all $k$, with $C$ depending only on $c$. From here on $C,C',c'$ denote
constants that depend only on $c$ and $D$, possibly different at each occurrence; the threshold $C$ of
Corollary~\ref{cor:rho} plays no further role.

\begin{lemma}\label{lem:folded}
\begin{enumerate}
\item[(i)] For each $s$, $j\mapsto p_s(j,0)$ is nonincreasing.
\item[(ii)] $s\mapsto p_s(0,0)$ is nonincreasing and $s\mapsto p_s(m,0)$ is nondecreasing; $\partial_sp_s(0,0)=p_s(1,0)-p_s(0,0)$ and $\partial_sp_s(m,0)=p_s(m-1,0)-p_s(m,0)$.
\item[(iii)] $V(k+e_a)\ge V(k)$ whenever $k_a<m$.
\item[(iv)] $V(k)=D\int_0^\infty\bigl(\prod_a\hat p_s(0)-\prod_a\hat p_s(k_a)\bigr)\,ds$.
\end{enumerate}
\end{lemma}

\begin{proof}
Two copies of $X$ started at $j<j'$ and run independently until they meet, and together afterwards, stay ordered,
because nearest-neighbour jumps of independent chains in continuous time do not occur simultaneously. This gives (i),
and also (iii): coupling coordinate $a$ from $k_a$ and $k_a+1$ in this way and the other coordinates identically, the
upper copy is at $0$ only when the lower one is, so the first time all coordinates are at $0$ is pathwise later from
$k+e_a$. For (ii), the formulas are the backward equations at $0$ and $m$; $p_s(0,0)/\pi(0)=\sum_i\varphi_i(0)^2e^{-\lambda_is}$
in an orthonormal eigenbasis of $L^2(\pi)$ is nonincreasing; and $X$ started at $m$ is stochastically nonincreasing in
time by the same coupling (compare the chain from $m$ with the chain from $X_r\le m$), so $P_m(X_s=0)$ is
nondecreasing. For (iv) let $\mathcal L$ be the generator of the product chain and
$Z(k)=\int_0^\infty(q_t(k)-\pi_Q(0))\,dt$ with $q_t(k)=\prod_ap_{t/D}(k_a,0)$; the integral converges because the product chain is finite and
irreducible, so $q_t(k)\to\pi_Q(0)$ exponentially fast. By the backward equation, $\mathcal LZ(k)=\int_0^\infty\partial_tq_t(k)\,dt=\pi_Q(0)-\mathbf
1\{k=0\}$, where $\mathbf 1$ denotes the indicator, so $W=(Z(0)-Z(\cdot))/\pi_Q(0)$ satisfies $\mathcal LW=-1$ off $0$ and $W(0)=0$. On a finite irreducible
chain this determines $W=V$. Substituting $t=Ds$ and~\eqref{eq:product} gives (iv).
\end{proof}

We need an exit bound and three kernel bounds for $X$.

\begin{lemma}\label{lem:exit}
For every $x$, $s>0$ and $r>0$ with $\eps s\le r/4$, $P_x\bigl(\sup_{u\le s}|X_u-x|\ge r\bigr)\le2\exp\bigl(-r^2/(5(s+r))\bigr)$.
\end{lemma}

\begin{proof}
Let $f(z)=\cosh(\theta(z-x))$ with $0<\theta\le1$. At an interior $z$,
$\alpha f(z+1)+\beta f(z-1)-f(z)=f(z)(\cosh\theta-1)+\eps\sinh(\theta(z-x))\sinh\theta\le\psi f(z)$ with
$\psi=\cosh\theta-1+\eps\sinh\theta$. At $z=0$ the increment is $f(1)-f(0)$, which is $\le0$ if $x\ge1$ and equals
$(\cosh\theta-1)f(0)$ if $x=0$; the same holds at $z=m$. So $e^{-\psi u}f(X_u)$ is a supermartingale. At the first time
$\tau$ with $|X_\tau-x|\ge r$ we have $f(X_\tau)\ge\cosh(\theta r)$, and optional stopping at $\tau\wedge s$ gives
$P(\tau\le s)\le2\exp(\psi s-\theta r)$. Since $\cosh\theta-1\le0.55\,\theta^2$ and $\sinh\theta\le1.18\,\theta$ for
$\theta\le1$, and $\eps s\le r/4$, we get $\psi s-\theta r\le0.55\,\theta^2s-0.7\,\theta r$. Take $\theta=r/(2s)$ if
$r\le2s$ and $\theta=1$ otherwise; in both cases the exponent is at most $-r^2/(5(s+r))$.
\end{proof}

\begin{lemma}\label{lem:kernel}
Let $\eps=c/L$ and $\mathcal S(s)=\max_x p_s(x,x)/\pi(x)$.
\begin{enumerate}
\item[(a)] $\pi(0)\,\mathcal S(s)\le C\bigl(1\wedge s^{-1/2}+L^{-1}\bigr)$; hence $\hat p_s(j)\le\mathcal S(s)\le C(L(1\wedge s^{-1/2})+1)$.
\item[(b)] $\hat p_s(j)\le\mathcal S(s/2)\bigl(P_j(X_{s/2}\le j/2)+P_0(X_{s/2}>j/2)\bigr)$; so if $\eps s\le j/4$, then
$\hat p_s(j)\le C(L(1\wedge s^{-1/2})+1)\exp\bigl(-j^2/(10(s+j))\bigr)$.
\item[(c)] $|\hat p_s(j)-1|\le Ce^{-(s-L^2)/(CL^2)}$ for $s\ge L^2$.
\end{enumerate}
\end{lemma}

\begin{proof}
(a) The chain $X$ is the cycle distance $\min(Y,L-Y)$ of the chain $Y$ on $\Z_L$ that steps away from $0$ at rate
$\alpha$ and toward $0$ at rate $\beta$, and from $0$ and from $L/2$ to each neighbour at rate $1/2$. Let $\tilde Y$
be its periodic lift to $\Z$ and $\tilde P$ the jump matrix of $\tilde Y$. With $r=\alpha/\beta$, the lazy matrix $(I+\tilde P)/2$ is
reversible with respect to the periodic measure $\tilde m$ given on one period by $\tilde m(0)=1$,
$\tilde m(i)=\tilde m(L-i)=r^{i-1}/(1-\eps)$ for $1\le i\le m-1$ and $\tilde m(m)=r^{m-1}$; the detailed-balance
equations $\tilde m(0)/2=\tilde m(1)\beta$, $\tilde m(i)\alpha=\tilde m(i+1)\beta$ and $\tilde m(m-1)\alpha=\tilde m(m)/2$
hold, and the other half of the period is their reflection. For $\eps\le1/2$ we have $\log r\le4\eps$,
so $1\le\tilde m\le e^{2c}$; the entries of $(I+\tilde P)/2$ are at least $1/8$ on edges and equal to $1/2$ on loops, so the
edge and loop weights $\tilde m(x)(I+\tilde P)(x,y)/2$ lie in $[1/8,e^{2c}]$. Balls in $\Z$ are intervals, so volume
doubling holds with constant $3e^{2c}$, and writing each difference $f(z)-f(x)$ as a sum along the interval and
applying the Cauchy--Schwarz inequality gives the Poincar\'e inequality on balls of radius $\rho$ with constant
$48e^{2c}\rho^2$. These are the hypotheses of Delmotte~\cite{Del99}, with constants depending only on $c$. Then \cite[Theorem~1.7]{Del99} gives
$((I+\tilde P)/2)^n(x,y)\le Cn^{-1/2}e^{-(x-y)^2/(Cn)}$ for $n\ge1$. Since
$p_s=\sum_ne^{-2s}(2s)^n/n!\,((I+P)/2)^n$ with $P$ the jump matrix of $X$, and $X_s=j$ means that $\tilde Y_s$ lies in
one of at most two residue classes mod $L$, summing over the images $y+zL$ gives $p_s(x,y)\le C\,\E[(1\wedge
N^{-1/2})+L^{-1}]$ with $N$ a Poisson variable of mean $2s$, which is at most $C(1\wedge s^{-1/2}+L^{-1})$. Together
with $\pi(y)\ge\pi(0)/C$ this gives (a). For the second part, by Cauchy--Schwarz in the eigenbasis,
$p_s(x,y)/\pi(y)\le\bigl(p_s(x,x)p_s(y,y)/(\pi(x)\pi(y))\bigr)^{1/2}\le\mathcal S(s)$.
(b) Split $p_s(j,0)=\sum_yp_{s/2}(j,y)p_{s/2}(y,0)$ at $y=j/2$. For $y\le j/2$ use $p_{s/2}(y,0)\le\pi(0)\mathcal
S(s/2)$; for $y>j/2$ use $p_{s/2}(j,y)\le\pi(y)\mathcal S(s/2)$ and $\pi(y)p_{s/2}(y,0)=\pi(0)p_{s/2}(0,y)$. The second
statement follows from Lemma~\ref{lem:exit} with $r=j/2$ and time $s/2$.
(c) The spectral gap $\lambda_1$ of $X$ is at least $1/(CL^2)$: the Dirichlet form and the variance change by
bounded factors under the comparison in (a), and the plain path has gap of order $m^{-2}$. In the eigenbasis,
$|\hat p_s(j)-1|\le\bigl((p_s(j,j)/\pi(j)-1)(p_s(0,0)/\pi(0)-1)\bigr)^{1/2}$ and
$p_s(x,x)/\pi(x)-1\le e^{-\lambda_1(s-L^2)}(\mathcal S(L^2)-1)$, which is at most $Ce^{-\lambda_1(s-L^2)}$ by (a).
\end{proof}

\section{A product inequality on the half-line}\label{sec:halfline}

In this section only, fix $\eps\in[0,1)$, put $\alpha=(1+\eps)/2$, $\beta=(1-\eps)/2$, and let $X^\infty$ be the
continuous-time chain on $\N=\{0,1,2,\dots\}$ that jumps from $0$ to $1$ at rate $1$ and, from $x\ge1$, to $x+1$ at
rate $\alpha$ and to $x-1$ at rate $\beta$. Write $f_t(x)=P_x(X^\infty_t=0)$.

\begin{lemma}\label{lem:halfline}
For all $t>0$ and $j\ge1$,
\[
  f_t(1)\,f_t(j)\ \ge\ f_t(0)\,f_t(j+1).
\]
\end{lemma}

For $\eps=0$ the chain is $|S|$ for the continuous-time simple random walk $S$ on $\Z$, so $f_t(j)=e^{-t}I_j(t)$ with
$I_j$ the modified Bessel functions, and the lemma says $I_{j+1}(t)/I_j(t)\le I_1(t)/I_0(t)$. This case is classical: it follows by chaining the
Tur\'an-type inequality $I_\nu(t)^2\ge I_{\nu-1}(t)I_{\nu+1}(t)$ over $\nu=1,\dots,j$; see~\cite{BP} for this inequality
and its history. Lemma~\ref{lem:halfline} extends it to $\eps>0$.

\begin{proof}[Proof of Lemma~\ref{lem:halfline}]
\emph{The pair chain.} Run two copies of $X^\infty$ from $j$ and $j+1$ as follows. While they differ they occupy
$(y,y+1)$; both move up at rate $\alpha$ and both move down at rate $\beta$ when $y\ge1$; when $y=0$ both move up at rate $\alpha$,
the lower one alone moves up at rate $\beta$ and the upper one alone moves down at rate $\beta$. Each copy has the law of
$X^\infty$, the copies stay ordered, and they move together once they meet. Hence
\[
  k_t(j):=f_t(j)-f_t(j+1)=P(\text{lower copy at }0\text{ at time }t,\ \text{copies not yet met}).
\]
The lower copy before the meeting is the chain $\tilde X$ on $\N$ with the interior rates of $X^\infty$, rate $\alpha$ from $0$ to
$1$, and killing at rate $2\beta$ at $0$; so $k_t(j)=P_j(\tilde X_t=0)$. Both $f$ and $k$ solve the backward equations of their
chains, $f_t,k_t>0$, and $0\le k_t\le f_t$.

\emph{The equation for the ratio.} Let $h_t=k_t/f_t\in[0,1]$. The two generators agree at $x\ge1$, and substituting
$k=hf$ gives
\begin{align*}
  \partial_t h_t(x)&=\alpha\frac{f_t(x+1)}{f_t(x)}\bigl(h_t(x+1)-h_t(x)\bigr)+\beta\frac{f_t(x-1)}{f_t(x)}\bigl(h_t(x-1)-h_t(x)\bigr),
  \qquad x\ge1,\\
  \partial_t h_t(0)&=\alpha\frac{f_t(1)}{f_t(0)}\bigl(h_t(1)-h_t(0)\bigr)-\kappa_t\,h_t(0),\qquad
  \kappa_t=\beta\Bigl(1+\frac{f_t(1)}{f_t(0)}\Bigr).
\end{align*}
There is no potential away from $0$. Fix $0<\eta<t$. For $s\in[0,t-\eta]$ let $\zeta$ be the time-inhomogeneous
nearest-neighbour chain that at time $s$ jumps up at rate $\alpha f_{t-s}(x+1)/f_{t-s}(x)\le\alpha$ and down at rate
$\beta f_{t-s}(x-1)/f_{t-s}(x)$ (here $f_{t-s}(x+1)\le f_{t-s}(x)$ by the monotone coupling). Its upward rates are bounded
and a downward jump cannot be repeated more than $x$ times from $x$, so $\zeta$ does not explode, and the number of its
upward jumps on $[0,t-\eta]$ is dominated by a Poisson variable of mean $\alpha t$. The process
$M_s=\exp\bigl(-\int_0^s\kappa_{t-r}\mathbf{1}\{\zeta_r=0\}\,dr\bigr)\,h_{t-s}(\zeta_s)$ is a bounded martingale, so
\[
  h_t(x)=\E_x\Bigl[\exp\Bigl(-\int_0^{t-\eta}\kappa_{t-r}\mathbf{1}\{\zeta_r=0\}\,dr\Bigr)\,h_\eta(\zeta_{t-\eta})\Bigr].
\]

\emph{Coupling.} Run $\zeta$ from $0$ and from $x$ independently until they meet and together afterwards. Then
$\zeta^0_s\le \zeta^x_s$ for all $s$; in particular $\zeta^x_s=0$ forces $\zeta^0_s=0$, so the killing weights satisfy $S^x\ge S^0$
pathwise, where $S^y$ denotes the exponential factor above for the copy started at $y$. Since $h_\eta\ge0$ and
$S^0\le1$, for every $N$,
\[
  h_t(x)-h_t(0)\ \ge\ \E\bigl[S^0\bigl(h_\eta(\zeta^x_{t-\eta})-h_\eta(\zeta^0_{t-\eta})\bigr)\bigr]
  \ \ge\ -2\max_{y\le N}|1-h_\eta(y)|-2P\bigl(\mathrm{Poisson}(\alpha t)>N-x\bigr).
\]
For fixed $y$, both $f_\eta(y)$ and $k_\eta(y)$ equal $(\beta\eta)^y/y!\,(1+O(\eta))$ as $\eta\to0$ (the only paths of
order $\eta^y$ are the $y$ straight downward steps, at rate $\beta$ in both chains), so $h_\eta(y)\to1$. Letting
$\eta\to0$ and then $N\to\infty$ gives $h_t(x)\ge h_t(0)$.

\emph{Conclusion.} $h_t(j)\ge h_t(0)$ reads $\bigl(f_t(j)-f_t(j+1)\bigr)f_t(0)\ge\bigl(f_t(0)-f_t(1)\bigr)f_t(j)$,
which is the claim.
\end{proof}

On the folded chain $X$ the same computation produces a second killing site, at $m-1$, with rate
$\alpha\bigl(1+p_t(m,0)/p_t(m-1,0)\bigr)$, and we do not use the inequality of Lemma~\ref{lem:halfline} for $X$
itself. Lemma~\ref{lem:Aprime} uses Lemma~\ref{lem:halfline} only at times
$t\le t_0=L^2/(K\log L)$, with the constant $K$ fixed in its proof, and for starts $j\le m/2$, where $X$ and $X^\infty$ coincide up to an event of small
probability; the case $j>m/2$ does not use it.

\section{The deficit}\label{sec:deficit}

Throughout, $\eps=c/L$. By Lemma~\ref{lem:folded}(iv), for any two sites $y,y'$,
\begin{equation}\label{eq:diff}
V(y)-V(y')=D\int_0^\infty\Bigl(\prod_a\hat p_s(y'_a)-\prod_a\hat p_s(y_a)\Bigr)\,ds .
\end{equation}

\begin{lemma}\label{lem:Aprime}
There is $C$ such that $V(k+e_a)-V(k+e_i)\le C(\log L)^{D/2}$ whenever $k_a=0$ and $1\le k_i\le m-1$.
\end{lemma}

\begin{proof}
Let $j=k_i$, $R(s)=\prod_{a'\ne a,i}\hat p_s(k_{a'})$, and
$\Gamma(s)=\hat p_s(1)\hat p_s(j)-\hat p_s(0)\hat p_s(j+1)$. By~\eqref{eq:diff},
$V(k+e_a)-V(k+e_i)=D\int_0^\infty R(s)(-\Gamma(s))\,ds\le D\int_0^\infty R(s)\Gamma(s)^-\,ds$. By
Lemma~\ref{lem:folded}(i), $0\le R(s)\le\hat p_s(0)^{D-2}$. Put $t_0=L^2/(K\log L)$ with $K$ chosen below.

\emph{Large times.} Since $\hat p_s(j+1)\le\hat p_s(j)\le\hat p_s(0)$ and $\hat p_s(0)-\hat p_s(1)=-\partial_s\hat p_s(0)$
(Lemma~\ref{lem:folded}(ii)),
\[
\Gamma(s)^-\le\hat p_s(j)\bigl(\hat p_s(0)-\hat p_s(1)\bigr)\le-\hat p_s(0)\,\partial_s\hat p_s(0).
\]
Hence, as $\hat p_s(0)\to1$,
\[
D\int_{t_0}^\infty R\Gamma^-\,ds\le-D\int_{t_0}^\infty\hat p_s(0)^{D-1}\partial_s\hat p_s(0)\,ds=\hat p_{t_0}(0)^D-1\le
C(K\log L)^{D/2}
\]
by Lemma~\ref{lem:kernel}(a).

\emph{Small times.} Let $X^\infty$ be the half-line chain of Section~\ref{sec:halfline} with the same $\eps$, coupled
with $X$ from the same start so that the two agree up to the hitting time $T_m$ of $m$. For $y\le j+1$,
$|p_s(y,0)-P_y(X^\infty_s=0)|\le P_y(T_m\le s)\le P_{j+1}(T_m\le s)$, the last step by the monotone coupling.

If $j\le m/2$, then $m-j-1\ge L/5$ for large $L$, and for $s\le t_0$ Lemma~\ref{lem:exit} (with $r=L/5$, which for large $L$ satisfies $r\le t_0$ and $\eps t_0\le r/4$,
so that $5(s+r)\le10t_0$) gives $P_{j+1}(T_m\le s)\le2\exp(-L^2/(250\,t_0))=2L^{-K/250}=:\eta_L$. Each of $p_s(0,0)$,
$p_s(1,0)$, $p_s(j,0)$, $p_s(j+1,0)$ lies in $[0,1]$ and is within $\eta_L$ of its half-line counterpart, and the
same bracket with $P_y(X^\infty_s=0)$ in place of $p_s(y,0)$ is nonnegative by Lemma~\ref{lem:halfline}, so $\pi(0)^2\Gamma(s)^-\le4\eta_L$.

If $j>m/2$, then $\pi(0)^2\Gamma(s)^-\le p_s(j+1,0)\le P_{j+1}(\sup_{u\le s}|X_u-X_0|\ge j+1)\le\eta_L$ in the same way.

In both cases $R(s)\Gamma(s)^-\le(CL)^{D-2}\cdot4\eta_L\,(CL)^2$ for $s\le t_0$, by Lemma~\ref{lem:kernel}(a) and
$\pi(0)^{-1}\le CL$. So $D\int_0^{t_0}R\Gamma^-\,ds\le C'L^{D+2}L^{-K/250}\le1$ once $K=250(D+3)$ and $L$ is large.
\end{proof}

\begin{lemma}\label{lem:B}
There is $C$ such that for every $k$ with $k_i=m$ for some $i$:
\begin{enumerate}
\item[(i)] $V(k)-V(k-e_i)\le C$;
\item[(ii)] $V(k+e_a)-V(k)\le C$ for every $a$ with $k_a=0$.
\end{enumerate}
\end{lemma}

\begin{proof}
(i) By~\eqref{eq:diff} and Lemma~\ref{lem:folded}(ii),
$V(k)-V(k-e_i)=D\int_0^\infty R_i(s)\,\partial_s\hat p_s(m)\,ds$ with $R_i(s)=\prod_{a\ne i}\hat p_s(k_a)\le\hat
p_s(0)^{D-1}$, and $\partial_s\hat p_s(m)\ge0$. On $[L^2,\infty)$, $R_i\le C$ by Lemma~\ref{lem:kernel}(a) and
$\int_{L^2}^\infty\partial_s\hat p_s(m)\,ds\le1$. Let $s_1=L^2/(8c+8)$. On $[s_1,L^2]$ we have
$R_i\le\hat p_{s_1}(0)^{D-1}\le C$ and $\int_{s_1}^{L^2}\partial_s\hat p_s(m)\,ds\le\hat p_{L^2}(m)\le C$. For $s\le
s_1$ we use that $s\mapsto\hat p_s(m)$ is nondecreasing with $\hat p_0(m)=0$, and that Lemma~\ref{lem:kernel}(b) with
$j=m$ applies at every $s\le s_1$ (since $\eps s_1\le m/4$) and gives
\begin{equation}\label{eq:pm}
\hat p_s(m)\le C\bigl(L(1\wedge s^{-1/2})+1\bigr)\exp\Bigl(-\frac{L^2}{C(s+L)}\Bigr),\qquad s\le s_1.
\end{equation}
On $(0,L]$, $R_i\le(CL)^{D-1}$ and $\int_0^L\partial_s\hat p_s(m)\,ds=\hat p_L(m)\le CLe^{-L/(2C)}$, so this range
contributes at most $(CL)^De^{-L/(2C)}$. Cover $[L,s_1]$ by the blocks $[U,2U]$, $U=2^{-r}s_1\ge L/2$, $r\ge1$. On such a block $R_i\le(C L/\sqrt U)^{D-1}$ and the
integral of $\partial_s\hat p_s(m)$ is at most $\hat p_{2U}(m)$, so
\[
\int_U^{2U}R_i\,\partial_s\hat p_s(m)\,ds\le\Bigl(\frac{C'L^2}{U}\Bigr)^{D/2}\exp\Bigl(-\frac{L^2}{4CU}\Bigr).
\]
Over these blocks the ratio $L^2/U$ grows geometrically from at least $8c+8$, so
the sum is bounded uniformly in $L$.

(ii) Similarly $V(k+e_a)-V(k)=D\int_0^\infty R_a(s)\bigl(-\partial_s\hat p_s(0)\bigr)\,ds$ with
$R_a(s)=\prod_{a'\ne a}\hat p_s(k_{a'})$, a product that contains the factor $\hat p_s(m)$ from coordinate $i$, and
$-\partial_s\hat p_s(0)\ge0$. On $[s_1,\infty)$, $R_a\le C$ and $\int_{s_1}^\infty(-\partial_s\hat p_s(0))\,ds\le\hat
p_{s_1}(0)\le C$. On $(0,L]$, $\hat p_s(m)\le\hat p_L(m)\le CLe^{-L/(2C)}$, the other $D-2$ factors of $R_a$ are at most
$CL$, and $\int_0^L(-\partial_s\hat p_s(0))\,ds\le\hat p_0(0)=\pi(0)^{-1}\le CL$, so this range contributes at most
$(CL)^De^{-L/(2C)}$. On each block $[U,2U]$ of (i), $\int_U^{2U}(-\partial_s\hat p_s(0))\,ds\le\hat p_U(0)\le
CL/\sqrt U$, $\hat p_s(m)\le\hat p_{2U}(m)$, and the other factors are at most $CL/\sqrt U$, which gives the block bound
of (i) and the same summation.
\end{proof}

\begin{proposition}\label{prop:deficit}
$\bar d=\max_{x\ne0}d(x)\le C(\log L)^{D/2}/L$.
\end{proposition}

\begin{proof}
If $S_m=\emptyset$, Lemma~\ref{lem:deficit} and Lemma~\ref{lem:Aprime} (applied with each $i\in I$) give
$d(x)\le(c/L)\,C(\log L)^{D/2}$ (if $I=\emptyset$ too, then $x=0$). If $S_m\ne\emptyset$, the second
bound of Lemma~\ref{lem:deficit} and Lemma~\ref{lem:B} give $d(x)\le(c/L)\,2C$.
\end{proof}

\section{The maximum expected hitting time under the product adversary}\label{sec:value}

\begin{proposition}\label{prop:value}
$L^{-D}\max_xV(x)\to G_D\bigl((e^c-1)/c\bigr)^D$ as $L\to\infty$, for every fixed $c\ge0$, with $(e^c-1)/c=1$ at $c=0$.
\end{proposition}

\begin{proof}
Let $q_t$ and $Z$ be as in the proof of Lemma~\ref{lem:folded}(iv), read as functions of $x$ through $k(x)$, so that
$q_t(x)=P^Q_x(\xi_t=0)$ by~\eqref{eq:product}, $V(x)=(Z(0)-Z(x))/\pi_Q(0)$ and $\max_xV=(Z(0)-\min_xZ(x))/\pi_Q(0)$. By~\eqref{eq:pi0},
$\pi_Q(0)L^D\to(c/(e^c-1))^D$, so it suffices to show $\min_xZ(x)\to0$ and $Z(0)\to G_D$.

\emph{The minimum.} Stationarity gives $\sum_x\pi_Q(x)Z(x)=0$, so $\min_xZ(x)\le0$. By~\eqref{eq:product} and
Lemma~\ref{lem:kernel}(c), $|q_t(x)-\pi_Q(0)|\le C\pi_Q(0)e^{-(t/D-L^2)/(CL^2)}$ for $t\ge2DL^2$, so
$Z(x)\ge-2DL^2\pi_Q(0)-CL^2\pi_Q(0)=-O(L^{2-D})$.

\emph{The value at $0$.} Let $q^\infty_t=P_0(S_t=0)$ for the continuous-time simple random walk $S$ on $\Z^D$ with jump
rate one, so that $\int_0^\infty q^\infty_t\,dt=G_D$. For fixed $t$, $q_t(0)\to q^\infty_t$: run $\xi$ and $S$ with the same
jump times and couple each jump direction maximally. The two laws of a jump differ by at most $\eps/2$ in total
variation, and as long as every coordinate of $S$ stays below $m$ in absolute value, $\xi_t=0$ if and only if $S_t=0$.
Hence $|q_t(0)-q^\infty_t|\le\eps t/2+P(N_t\ge m)$ with $N_t$ a Poisson variable of mean $t$, which tends to $0$. For
$t\le L^2$, Lemma~\ref{lem:kernel}(a) and~\eqref{eq:product} give
$q_t(0)=p_{t/D}(0,0)^D\le\bigl(C(1\wedge(t/D)^{-1/2}+L^{-1})\bigr)^D\le C'(1\wedge t^{-1/2})^D$, since $L^{-1}\le1\wedge(t/D)^{-1/2}$
there. The bound is integrable on $(0,\infty)$ because $D\ge3$, so by dominated convergence
$\int_0^{L^2}q_t(0)\,dt\to G_D$. On $[L^2,2DL^2]$ the same bound gives $q_t(0)\le CL^{-D}$, the range $t\ge2DL^2$
contributes $O(L^{2-D})$ as in the previous paragraph, and $\int_0^{2DL^2}\pi_Q(0)\,dt=O(L^{2-D})$. Hence
$Z(0)=\int_0^{L^2}q_t(0)\,dt+O(L^{2-D})\to G_D$.
\end{proof}

\begin{proof}[Proof of Theorem~\ref{thm:main}]
By Proposition~\ref{prop:deficit}, $\bar d=O((\log L)^{D/2}/L)<1$ for large $L$, so Lemma~\ref{lem:reduction} gives
$\thit^+(c/L)=\max_xV(x)\,(1+O((\log L)^{D/2}/L))$. Proposition~\ref{prop:value} for $c$ and for $c=0$ gives
$L^{-D}\max_xV(x)\to G_D((e^c-1)/c)^D$ and $L^{-D}\thit\to G_D$, and the ratio of the two limits is $((e^c-1)/c)^D$.
\end{proof}

\section{Numerical values for $D=2$ and $D=3$}\label{sec:checks}

The exact value of $\thit^+(\eps)$ on a finite torus is the solution of a Markov decision problem: by
the argument of Lemma~\ref{lem:reduction}, applied to the solution $V^*$ of $V^*=1+BV^*$ off the target (with
$V^*(0)=0$), the history-independent adversary that attains the maximum in $BV^*$ at every site attains
$\thit^+(\eps)$. We computed $V^*$ by policy iteration on $\Z_L^2$ for $L\le128$, and $\max_xV_Q(x)$ on $\Z_L^3$ for
$L\le40$. For $D=2$ and $\eps=1/L$ the results are as follows.
\begin{center}
\begin{tabular}{rccccc}
\toprule
$L$ & 8 & 16 & 32 & 64 & 128\\ \midrule
$\thit$ & 111.9 & 558.8 & 2685.5 & 12547.9 & 57419.8\\
$\thit^+(1/L)/\thit$ & 2.0804 & 2.2674 & 2.3975 & 2.4860 & 2.5489\\ \bottomrule
\end{tabular}
\end{center}
The differences from $(e-1)^2=2.9525$, multiplied by $\log L$, are $1.81,1.90,1.92,1.94,1.96$, still increasing at
$L=128$; Theorem~\ref{thm:main} does not cover $D=2$, where $G_2=\infty$ and
$\thit$ is of order $L^2\log L$. For $D=3$ and $c=1$, $L\,(5.0732-\max_xV_Q(x)/\thit)$ equals $18.27,20.22,21.60,22.67$ at
$L=16,24,32,40$, consistent with a correction of order $(\log L)/L$; here $5.0732=(e-1)^3$. The deficit $\bar d$ of
Section~\ref{sec:reduction} computed numerically is $0.414$ and $0.266$ at $L=16,24$ ($D=3$, $c=1$), attained at sites
with one coordinate equal to $m$ and the others $0$. The code is in the repository~\cite{Code}.

On the box $\{1,\dots,L\}^D$ the walk is not transitive and not regular at the boundary, and the order of $\rhit$
there is open.

\section*{Acknowledgements and funding}
This work was done by AI agents, Claude (Anthropic), connected to Math Brain and working inside the Brain framework, a
shared, persistent memory in which many agents, Claude and Codex (OpenAI), record claims, proofs, failed attempts
and methods, retrieve one another's results, coordinate who works on what, and audit each other's work. For more on the Brain and Math Brain,
see the Math Brain paper~\cite{MathBrainPaper}, the Math Brain repository~\cite{MathBrainRepo} and the Brain framework
repository~\cite{BrainFramework}.

The research presented in this paper was supported by the European Research Council (ERC) under the European Union's
Horizon Europe research and innovation programme (grant agreement No.~101041711), by the Simons Foundation as part of
the Collaboration on the Mathematical and Scientific Foundations of Deep Learning, by Heights Labs, by Convex Nexus
Capital, by the Israel Science Foundation (grant number 2258/19), and by the Israel Science Foundation (ISF Grant
4101/25).

\section*{Declaration of competing interest}
The author holds equity in Heights Labs and in Convex Nexus Capital and has a management role at Convex Nexus Capital.
These entities are included in the funding declaration above.

\section*{Data availability}
The programs described in Section~\ref{sec:checks} are in the repository~\cite{Code}.

\section*{Declaration of generative AI and AI-assisted technologies}
This paper is part of a broader research project on AI proving problems. Claude agents (Anthropic) proved
Theorem~\ref{thm:main}, wrote the code of Section~\ref{sec:checks} and prepared the manuscript, and a separate
Claude model reviewed the proofs and proposed the estimates of Section~\ref{sec:deficit}. The model, the
definitions and the question are due to Zuckerman~\cite{Zuc26}, the upper bound $\rhit=O(1/L)$ follows from the
adversary constructed in the proof of his Theorem~4.5, and the on-diagonal heat-kernel bound of
Lemma~\ref{lem:kernel}(a) rests on Delmotte's Gaussian estimates~\cite{Del99}; the case $\eps=0$ of Lemma~\ref{lem:halfline} is a Tur\'an-type
inequality for modified Bessel functions~\cite{BP}; Lemma~\ref{lem:reduction} is a standard supermartingale comparison. The author directed the work and is responsible
for the content of the paper and for decisions concerning its dissemination.

\appendix
\section{Notation}\label{app:notation}
\begin{center}
\begin{tabular}{@{}lp{0.5\textwidth}l@{}}
\toprule
Symbol & Meaning & Where\\ \midrule
$G$, $P$ & a connected graph and its simple random walk & Section~\ref{sec:intro}\\
$\eps$, $c$ & adversary strength, $\eps=c/L$ & Section~\ref{sec:intro}\\
$\thit$, $\thit^+(\eps)$ & maximum expected hitting time, plain and under standard $\eps$-adversaries & Section~\ref{sec:intro}\\
$\rhit(C)$ & hitting-time robustness at threshold $C$ & Section~\ref{sec:intro}\\
$H_v$, $h^+(u,v)$ & hitting time of $v$; its supremum expectation from $u$ over standard adversaries & Section~\ref{sec:intro}\\
$\Delta$, $\bar\Delta$, $t_{\mathrm{mix}}$ & diameter, average distance, mixing time & Section~\ref{sec:intro}\\
$L$, $D$, $m$ & side, dimension, $m=L/2$ & Section~\ref{sec:intro}\\
$k=(k_a)$ & folded coordinates, $k_a=\min(x_a,L-x_a)$; $e_a$ the $a$-th unit vector & Section~\ref{sec:intro}\\
$Q$, $V=V_Q$ & the product adversary and $V(x)=\E^Q_xH_0$ & Section~\ref{sec:intro}\\
$G_D$, $r_D(C)$ & Green constant of $\Z^D$ at $0$; the limit of $L\rhit(C)$ & Theorem~\ref{thm:main}, Corollary~\ref{cor:rho}\\
$B$, $d$, $\bar d$ & best one-step operator, deficit $d=BV-QV$, its maximum & Section~\ref{sec:reduction}\\
$\alpha$, $\beta$ & rates $(1\pm\eps)/2$ away from and toward $0$ & Section~\ref{sec:folded}\\
$X$, $\pi$, $p_s$ & folded chain, its stationary law, its transition probabilities & Section~\ref{sec:folded}\\
$\hat p_s(j)$, $\mathcal S(s)$ & $p_s(j,0)/\pi(0)$; $\max_xp_s(x,x)/\pi(x)$ & Section~\ref{sec:folded}\\
$\pi_Q$, $Z$ & stationary law of the product chain; the integrated kernel defining $V$ & Section~\ref{sec:folded}\\
$X^\infty$, $f_t$ & the half-line chain; $f_t(x)=P_x(X^\infty_t=0)$ & Section~\ref{sec:intro}, Section~\ref{sec:halfline}\\
$h_t$, $\kappa_t$, $\zeta$ & the ratio $k_t/f_t$ with $k_t(j)=f_t(j)-f_t(j+1)$, the killing rate at $0$, the time-inhomogeneous chain & Section~\ref{sec:halfline}\\
$\Gamma(s)$ & $\hat p_s(1)\hat p_s(j)-\hat p_s(0)\hat p_s(j+1)$ & Section~\ref{sec:deficit}\\
$t_0$, $K$, $\eta_L$, $T_m$ & $L^2/(K\log L)$, $K=250(D+3)$, $2L^{-K/250}$, the hitting time of $m$ & Section~\ref{sec:deficit}\\
\bottomrule
\end{tabular}
\end{center}

\begin{thebibliography}{9}
\bibitem{Anon26} Anonymous, \emph{Small diameter does not control hitting-time robustness}, preprint, Zenodo, 2026, \doi{10.5281/zenodo.23575572}.
\bibitem{BP} \'A.~Baricz and S.~Ponnusamy, \emph{On Tur\'an type inequalities for modified Bessel functions},
Proc.\ Amer.\ Math.\ Soc.\ \textbf{141} (2013), 523--532, \doi{10.1090/S0002-9939-2012-11325-5}.
\bibitem{BrainFramework} Y.~Shmalo, \emph{Brain Framework}, public software repository, 2026,
\url{https://github.com/yspennstate/brain-framework-public/tree/5629e644cdd510173ae97e9aa92642dabd22dc57}.
\bibitem{Code} Y.~Shmalo, \emph{Hitting-time robustness of the discrete torus: source code and paper}, 2026,
\url{https://github.com/yspennstate/torus-hitting-robustness}.
\bibitem{Del99} T.~Delmotte, \emph{Parabolic Harnack inequality and estimates of Markov chains on graphs}, Rev.\ Mat.\
Iberoamericana \textbf{15} (1999), 181--232, \doi{10.4171/RMI/254}.
\bibitem{LPW17} D.~A.~Levin and Y.~Peres, \emph{Markov chains and mixing times}, 2nd ed., with contributions by
E.~L.~Wilmer, American Mathematical Society, Providence, RI, 2017, \doi{10.1090/mbk/107}.
\bibitem{MathBrainRepo} Y.~Shmalo, \emph{Math Brain: public snapshot, code and tests}, 2026,
\url{https://github.com/yspennstate/math-brain-public/tree/8c9a6bd478c4ec4082347151b7ae16db50c917dd}.
\bibitem{MathBrainPaper} Y.~Shmalo, \emph{Math Brain: shared, auditable memory for multi-agent mathematical research},
2026, \doi{10.5281/zenodo.24150837},
\url{https://github.com/yspennstate/math-brain-public/tree/main/papers/math-brain-2026-10-24}.
\bibitem{ShmCube} Y.~Shmalo, \emph{Hitting-time robustness of the hypercube}, preprint, 2026,
\url{https://github.com/yspennstate/hypercube-hitting-robustness}.
\bibitem{Zuc26} D.~Zuckerman, \emph{Markov chain robustness}, in: 17th Innovations in Theoretical Computer Science
Conference (ITCS 2026), LIPIcs~362, Schloss Dagstuhl--Leibniz-Zentrum f\"ur Informatik, 2026, Paper~118,
\doi{10.4230/LIPIcs.ITCS.2026.118}.
\end{thebibliography}

\end{document}
